\documentclass[main.tex]{subfiles}
\begin{document}
% ============================================================
%  8. Globalization: line search along the projection arc   (~10 min)
% ============================================================
\section{Line Search Along the Projection Arc}\label{sec:linesearch}

\begin{frame}{The projection arc}
  \[
    u(s) := \Proj\big(u - s\nabla f(u)\big), \qquad s \ge 0 .
  \]
  \begin{itemize}
    \item Searching along a \emph{fixed} direction $d$ with $u + sd \in \Uad$ forces tiny steps near the boundary.
    \item The arc $s \mapsto u(s)$ bends along the boundary of $\Uad$: \alert{many constraints can become active in one step}.
    \item Key property: $\ip{\nabla f(u)}{u - u(s)} \ge \tfrac1s\norm{u - u(s)}^2$.
  \end{itemize}
  \todo{TikZ picture: box, iterate, straight ray vs.\ bent projection arc}
\end{frame}

\begin{frame}{Armijo rule along the arc}
  \framesubtitle{Goldstein (1964), Levitin--Polyak (1966), Bertsekas (1976)}
  Given $\sigma \in (0,1)$, $\beta \in (0,1)$, $\bar s > 0$: choose $s_k = \beta^{m_k}\bar s$ with $m_k$ the smallest integer such that
  \[
    f(u_k) - f\big(u_k(s_k)\big) \;\ge\; \sigma\,\ip{\nabla f(u_k)}{u_k - u_k(s_k)}_U .
  \]
  Variant (e.g.\ Kelley 1999):
  \[
    f\big(u_k(s_k)\big) - f(u_k) \;\le\; -\frac{\sigma}{s_k}\norm{u_k - u_k(s_k)}_U^2 .
  \]
  \begin{itemize}
    \item Each backtracking step costs one projection and one state solve, \emph{no} new adjoint.
    \item Both conditions are stated in the $U$ inner product: consistent across meshes.
  \end{itemize}
\end{frame}

\begin{frame}{Convergence theory}
  \begin{theorem}[Bertsekas (1976); Dunn (1981) in Hilbert space]
    Let $\nabla f$ be Lipschitz on $\Uad$. Then the Armijo step is well defined, and every accumulation point of the projected gradient iterates is stationary.
  \end{theorem}
  \textbf{Dunn's contributions}
  \begin{itemize}
    \item Global convergence and rate estimates for projected gradient processes in Hilbert space.
    \item Role of nondegeneracy/growth conditions at $\ub$ for linear rates.
    \item \alert{Active set identification}: under strict complementarity the active set is identified after finitely many steps (finite dim.) \todo{state precisely the function-space version you want to use}.
  \end{itemize}
  \vspace{0.5ex}
  This motivates: once the active set is known, switch to Newton on the free part.
\end{frame}

\end{document}
